\documentclass[../examen.tex]{subfiles}

\usepackage{amsmath}

\begin{document}

\noindent
\begin{tabular}{ll}
\textbf{Evaluación:} & Escuela Técnica N° 4-117 Ejército de los Andes - Termodinámica y Máquinas Térmicas 5°2°" \\
 & \\

\textbf{Nombre y Apellido:} & \hrulefill \\
\end{tabular}

\begin{enumerate}

    \item \textbf{(1.25 puntos)} Usando la ecuación de estado de los gases ideales, \( PV = nRT \), calcula la presión de 1 mol de gas ideal contenido en un volumen de 22.4 L a una temperatura de 273 K.

    \item \textbf{(1.25 puntos)} En un proceso isobárico, el volumen de un gas cambia de 1 m³ a 3 m³ mientras la temperatura se duplica. ¿Cuál es el trabajo realizado por el gas si la presión es constante a 1 atm?

    \item \textbf{(1.25 puntos)} Un gas ideal está a una presión de 3 atm y un volumen de 10 L. Si el gas se expande a 30 L y la temperatura permanece constante, ¿cuál es la nueva presión según la ley de Boyle-Mariotte?
    
    \item \textbf{(1.25 puntos)} Un gas ideal está contenido en un recipiente con volumen inicial de 2 L a una temperatura de 300 K. Según la ley de Avogadro, ¿qué ocurrirá si se duplica la cantidad de gas en el recipiente manteniendo constante la temperatura y el volumen?

    \item \textbf{(1.25 puntos)} Un cilindro con un pistón móvil contiene gas a una temperatura de 273 K y un volumen de 1 m³. De acuerdo a la ley de Charles, ¿qué volumen ocupará el gas si se calienta a 546 K, manteniendo la presión constante?

    \item \textbf{(1.25 puntos)} Un gas se encuentra en un contenedor de volumen constante. Según la ley de Gay-Lussac, ¿cómo cambia la presión del gas si su temperatura aumenta de 300 K a 600 K?

    \item \textbf{(1.25 puntos)} Un gas tiene un volumen de 5 L y una presión de 2 atm. Según la ley de Boyle-Mariotte, ¿cuál será el nuevo volumen si la presión se reduce a 1 atm mientras la temperatura se mantiene constante?

    

    \item \textbf{(1.25 puntos)} Según el primer principio de la termodinámica, si un gas recibe 500 J de calor y realiza un trabajo de 200 J, ¿cuál es el cambio en la energía interna del gas?
\end{enumerate}

\vspace{0.6cm}
\rule{.9\textwidth}{.1mm}
\vspace{0.6cm}

\end{document}
